\QHeading{第1题　口算。}
\par\noindent


\begin{oralgrid}
\OralItem{1}{125\times40}{5000}
\OralItem{2}{120\div30}{4}
\OralItem{3}{360\div60}{6}
\OralItem{4}{45\times20}{900}
\OralItem{5}{720\div80}{9}
\OralItem{6}{18\times50}{900}
\OralItem{7}{630\div90}{7}
\OralItem{8}{20\times35}{700}
\OralItem{9}{480\div16}{30}
\OralItem{10}{125\times0}{0}
\OralItem{11}{250\times4}{1000}
\OralItem{12}{24\times50}{1200}\end{oralgrid}
\par

\QHeading{口算思路补充（对应上表题号）}

\Answer{上面列出得数，下面按小题给出不同计算思路。计数单位法适用于小数，整数题可比较凑整、拆分或商不变等方法。}
\MethodTitle{第1小题：$125\times 40$}

\MethodBlock{口算思路一：拆因数，凑整计算}{把一个因数拆成两个因数，先凑成整百或整千。}{\begin{aligned}[t]&\phantom{={}}125\times 40\\
&=125\times (8\times 5)\\
&=(125\times 8)\times 5\\
&=1000\times 5\\
&=5000\end{aligned}}

\MethodBlock{口算思路二：按数位拆分计算}{把一个因数拆成整十、整百数与余下部分之和。}{\begin{aligned}[t]&\phantom{={}}125\times 40\\
&=40\times (100+25)\\
&=40\times 100+40\times 25\\
&=4000+1000\\
&=5000\end{aligned}}
\MethodTitle{第2小题：$120\div 30$}

\MethodBlock{口算思路一：同时去掉相同个数的0}{}{\begin{aligned}[t]&\phantom{={}}120\div 30\\
&=(120\div 10)\div (30\div 10)\\
&=12\div 3\\
&=4\end{aligned}}

\MethodBlock{口算思路二：拆分除数，连续相除}{}{\begin{aligned}[t]&\phantom{={}}120\div 30\\
&=120\div (2\times 15)\\
&=120\div 2\div 15\\
&=60\div 15\\
&=4\end{aligned}}
\MethodTitle{第3小题：$360\div 60$}

\MethodBlock{口算思路一：同时去掉相同个数的0}{}{\begin{aligned}[t]&\phantom{={}}360\div 60\\
&=(360\div 10)\div (60\div 10)\\
&=36\div 6\\
&=6\end{aligned}}

\MethodBlock{口算思路二：拆分除数，连续相除}{}{\begin{aligned}[t]&\phantom{={}}360\div 60\\
&=360\div (2\times 30)\\
&=360\div 2\div 30\\
&=180\div 30\\
&=6\end{aligned}}
\MethodTitle{第4小题：$45\times 20$}

\MethodBlock{口算思路一：按数位拆分计算}{把一个因数拆成整十、整百数与余下部分之和。}{\begin{aligned}[t]&\phantom{={}}45\times 20\\
&=20\times (40+5)\\
&=20\times 40+20\times 5\\
&=800+100\\
&=900\end{aligned}}
\MethodTitle{第5小题：$720\div 80$}

\MethodBlock{口算思路一：同时去掉相同个数的0}{}{\begin{aligned}[t]&\phantom{={}}720\div 80\\
&=(720\div 10)\div (80\div 10)\\
&=72\div 8\\
&=9\end{aligned}}

\MethodBlock{口算思路二：拆分除数，连续相除}{}{\begin{aligned}[t]&\phantom{={}}720\div 80\\
&=720\div (2\times 40)\\
&=720\div 2\div 40\\
&=360\div 40\\
&=9\end{aligned}}
\MethodTitle{第6小题：$18\times 50$}

\MethodBlock{口算思路一：拆因数，凑整计算}{把一个因数拆成两个因数，先凑成整百或整千。}{\begin{aligned}[t]&\phantom{={}}18\times 50\\
&=50\times (2\times 9)\\
&=(50\times 2)\times 9\\
&=100\times 9\\
&=900\end{aligned}}

\MethodBlock{口算思路二：凑整后补偿}{把接近整十、整百的因数凑整，再补回或减去多算的部分。}{\begin{aligned}[t]&\phantom{={}}18\times 50\\
&=50\times (20-2)\\
&=50\times 20-50\times 2\\
&=1000-100\\
&=900\end{aligned}}
\MethodTitle{第7小题：$630\div 90$}

\MethodBlock{口算思路一：同时去掉相同个数的0}{}{\begin{aligned}[t]&\phantom{={}}630\div 90\\
&=(630\div 10)\div (90\div 10)\\
&=63\div 9\\
&=7\end{aligned}}

\MethodBlock{口算思路二：拆分除数，连续相除}{}{\begin{aligned}[t]&\phantom{={}}630\div 90\\
&=630\div (2\times 45)\\
&=630\div 2\div 45\\
&=315\div 45\\
&=7\end{aligned}}
\MethodTitle{第8小题：$20\times 35$}

\MethodBlock{口算思路一：按数位拆分计算}{把一个因数拆成整十、整百数与余下部分之和。}{\begin{aligned}[t]&\phantom{={}}20\times 35\\
&=20\times (30+5)\\
&=20\times 30+20\times 5\\
&=600+100\\
&=700\end{aligned}}
\MethodTitle{第9小题：$480\div 16$}

\MethodBlock{口算思路一：利用商不变的性质}{被除数和除数同时乘同一个非零数，商不变。}{\begin{aligned}[t]&\phantom{={}}480\div 16\\
&=(480\times 5)\div (16\times 5)\\
&=2400\div 80\\
&=30\end{aligned}}

\MethodBlock{口算思路二：拆分除数，连续相除}{}{\begin{aligned}[t]&\phantom{={}}480\div 16\\
&=480\div (2\times 8)\\
&=480\div 2\div 8\\
&=240\div 8\\
&=30\end{aligned}}
\MethodTitle{第10小题：$125\times 0$}

\MethodBlock{口算思路一：按数位拆分计算}{把一个因数拆成整十、整百数与余下部分之和。}{\begin{aligned}[t]&\phantom{={}}125\times 0\\
&=0\times (100+25)\\
&=0\times 100+0\times 25\\
&=0+0\\
&=0\end{aligned}}
\MethodTitle{第12小题：$24\times 50$}

\MethodBlock{口算思路一：拆因数，凑整计算}{把一个因数拆成两个因数，先凑成整百或整千。}{\begin{aligned}[t]&\phantom{={}}24\times 50\\
&=50\times (2\times 12)\\
&=(50\times 2)\times 12\\
&=100\times 12\\
&=1200\end{aligned}}

\MethodBlock{口算思路二：按数位拆分计算}{把一个因数拆成整十、整百数与余下部分之和。}{\begin{aligned}[t]&\phantom{={}}24\times 50\\
&=50\times (20+4)\\
&=50\times 20+50\times 4\\
&=1000+200\\
&=1200\end{aligned}}


\par\noindent

\QHeading{第2题　脱式计算。}

\CalcBlock{1}{\begin{aligned}[t]&\phantom{={}}\color{black}76\times[288\div(130-106)]\\&\color{answerred}=76\times[288\div24]\\&\color{answerred}=76\times12\\&\color{answerred}=912\end{aligned}}
\par

\CalcBlock{2}{\begin{aligned}[t]&\phantom{={}}\color{black}592\div(61-24)\times26\\&\color{answerred}=592\div37\times26\\&\color{answerred}=16\times26\\&\color{answerred}=416\end{aligned}}
\CalcBlock{3}{\begin{aligned}[t]&\phantom{={}}\color{black}110\times[125\div(43-18)]\\&\color{answerred}=110\times[125\div25]\\&\color{answerred}=110\times5\\&\color{answerred}=550\end{aligned}}
\CalcBlock{4}{\begin{aligned}[t]&\phantom{={}}\color{black}16\times[(160+200)\div60]\\&\color{answerred}=16\times[360\div60]\\&\color{answerred}=16\times6\\&\color{answerred}=96\end{aligned}}
\CalcBlock[方法一]{5}{\begin{aligned}[t]&\phantom{={}}\color{black}864\div[(43-37)\times6]\\&\color{answerred}=864\div[6\times6]\\&\color{answerred}=864\div36\\&\color{answerred}=24\end{aligned}}
\MethodBlock{方法二：除以积，改成连续除}{连续除以两个数，等于除以这两个数的积。}{\begin{aligned}[t]&\phantom{={}}864\div [(43-37)\times 6]\\
&=864\div (43-37)\div 6\\
&=864\div 6\div 6\\
&=144\div 6\\
&=24\end{aligned}}

\CalcBlock{6}{\begin{aligned}[t]&\phantom{={}}\color{black}125\times[78\div(41-15)]\\&\color{answerred}=125\times[78\div26]\\&\color{answerred}=125\times3\\&\color{answerred}=375\end{aligned}}
\par\noindent

\QHeading{第3题　解决问题。}

\Problem{①修路队计划每天修路72米，增添设备后实际每天比计划多修18米。原来计划15天修完的路，现在多少天就可以修完？}
\par


\MethodTitle{方法一}
\Answer{\Formula{\begin{aligned}[t]&\phantom{={}}72\times15\div(72+18)\\&=72\times15\div90\\&=1080\div90\\&=12\text{（天）}\end{aligned}}\par 答：现在12天可以修完。}

\MethodTitle{方法二：先求可以少修几天}
\Answer{实际每天多修18米，若仍修15天就多修270米；用实际每天的修路量换算成少修的天数。}
\Answer{\textbf{多修的路程}\par \Formula{\begin{aligned}[t]&\phantom{={}}18\times 15\\
&=270\text{（米）}\end{aligned}}}
\Answer{\textbf{实际每天修路}\par \Formula{\begin{aligned}[t]&\phantom{={}}72+18\\
&=90\text{（米）}\end{aligned}}}
\Answer{\textbf{可以少修的天数}\par \Formula{\begin{aligned}[t]&\phantom{={}}270\div 90\\
&=3\text{（天）}\end{aligned}}}
\Answer{\textbf{实际所需天数}\par \Formula{\begin{aligned}[t]&\phantom{={}}15-3\\
&=12\text{（天）}\end{aligned}}}
\Answer{答：现在12天就可以修完。}
\Problem{②学校组织四年级的280名少先队员到敬老院参加活动，把这些少先队员平均分成7队，每队4组。平均每组有多少名少先队员？}

\MethodTitle{方法一}
\Answer{\Formula{\begin{aligned}[t]&\phantom{={}}280\div7\div4\\&=40\div4\\&=10\text{（名）}\end{aligned}}\par 答：平均每组有10名少先队员。}

\MethodTitle{方法二：先求一共有多少组}
\Answer{每队4组，共7队，先求总组数。}
\Answer{\Formula{\begin{aligned}[t]&\phantom{={}}280\div (4\times 7)\\
&=280\div 28\\
&=10\text{（名）}\end{aligned}}}
\Answer{答：平均每组有10名少先队员。}
